\documentclass[11pt,a4paper]{article}

\usepackage[utf8]{inputenc}
\usepackage[T1]{fontenc}
\usepackage{lmodern}
\usepackage[margin=1in]{geometry}
\usepackage{graphicx}
\usepackage{amsmath}
\usepackage{booktabs}
\usepackage{hyperref}
\usepackage{url}
\usepackage{caption}
\usepackage{subcaption}
\usepackage{multirow}
\usepackage{array}
\usepackage{parskip}
\usepackage{natbib}
\usepackage{float}
\usepackage{xcolor}

\hypersetup{
    colorlinks=true,
    linkcolor=blue,
    citecolor=blue,
    urlcolor=blue
}

\setlength{\parskip}{6pt}
\setlength{\parindent}{0pt}

\title{\textbf{How Robust Are Intent Classifiers to ASR Noise?\\
A Systematic Comparison of Traditional Models, BERT, and Large Language Models}}

\author{
    Abdul Moiz Muhammad\textsuperscript{1} \and
    Abdul Wasae\textsuperscript{1}
    \\[0.5em]
    \textsuperscript{1}Department of Computer Science, COMSATS University Islamabad\\
    \texttt{mojuaries111@gmail.com} \quad \texttt{abdulwasae1@gmail.com}
}

\date{May 2026}

\begin{document}

\maketitle

\begin{abstract}
Voice-based systems depend on a pipeline where automatic speech recognition (ASR) transcribes audio into text, which is then classified by an intent model. ASR errors degrade downstream classification accuracy, but no published work has compared how three major model families - TF-IDF classifiers, BERT-based encoders, and instruction-tuned large language models (LLMs) - degrade across a systematic range of word error rates on the same benchmark. We address this gap using SLURP \citep{bastianelli2020slurp}, a large-scale spoken language benchmark covering 66 intent classes. We fine-tune a Qwen2.5-3B model using QLoRA and evaluate all models at ASR word error rates (WER) from 0\% to 50\%. We find that the fine-tuned LLM maintains the highest Macro F1 at every noise level above zero. At WER=50\%, the LLM achieves a Macro F1 of 0.5888 compared to 0.5144 for BERT, a gap of 7.4 points. More notably, BERT degrades by 28.2 percentage points from clean to WER=50\%, while the LLM degrades by only 20.8 points. TF-IDF models show a different degradation pattern: a sharp initial drop at WER=10\% followed by relative stability up to WER=30\%, suggesting their sensitivity is front-loaded at low noise levels. We also establish the first published classification baselines on SLURP using modern LLMs, and release all noisy test sets for reproducibility.
\end{abstract}

\textbf{Keywords:} spoken language understanding, intent classification, ASR robustness, large language models, QLoRA, word error rate, SLURP

\section{Introduction}

Intent classification from transcribed speech is a core task in voice-based applications. A customer asking a smart speaker to set an alarm, a caller navigating an automated phone menu, or a user requesting information from a voice assistant all depend on systems that first transcribe speech and then classify the intent of the resulting text. The performance of the classification component is therefore constrained not only by the quality of the model itself, but also by the quality of the upstream ASR transcription.

Commercial ASR systems have improved significantly in recent years, but they remain imperfect. Under ideal conditions with clear speech from native speakers, modern ASR achieves word error rates below 5\%. Under realistic deployment conditions - accented speech, background noise, telephone-quality audio, or domain-specific vocabulary - WER commonly rises to 20-40\% or higher \citep{watanabe2017hybrid}. These errors are not random; they tend to cluster on semantically important words (proper nouns, domain-specific terms, numbers) that are precisely the words most useful for intent classification. The result is a systematic degradation of downstream classification that is difficult to characterize without controlled experiments.

Prior work has studied noise effects on individual model families. Several papers have shown that BERT's accuracy drops substantially as WER increases \citep{faruqui2021revisiting}. Work on the Audio-Snips dataset showed BERT accuracy falling to 75.86\% when WER reaches 52.94\% \citep{audio_snips_paper}. However, these studies evaluate a single model class in isolation. No published paper has placed TF-IDF classifiers, BERT-family encoders, and modern instruction-tuned LLMs on the same benchmark under the same noise conditions with the same evaluation protocol. Without this comparison, practitioners cannot make an informed choice between model architectures for noisy deployment environments.

This paper addresses that gap directly. Our contributions are as follows. First, we establish the first classification benchmarks on the SLURP dataset \citep{bastianelli2020slurp} using modern LLMs, providing a baseline that future work can compare against. Second, we present a controlled comparison of four model configurations across six WER levels (0, 10, 20, 30, 40, 50\%), identifying which architecture degrades most and least gracefully. Third, we identify and explain an unexpected pattern in TF-IDF degradation that differs qualitatively from transformer-based models. Fourth, we release all synthetically noisy test sets as CSV files to enable full reproducibility without requiring actual ASR systems.

\section{Related Work}

\subsection{Intent classification}

Intent classification has a long history in natural language processing, dating to rule-based and keyword-matching systems in the 1990s \citep{tur2011spoken}. Modern approaches generally fall into two categories: discriminative classifiers trained on text features, and fine-tuned pretrained language models.

In the discriminative category, TF-IDF combined with linear classifiers such as SVM or logistic regression remains a strong baseline on many intent datasets \citep{joulin2017bag}. These models are fast, interpretable, and require little memory at inference time, making them attractive for resource-constrained deployments. Their weakness is sensitivity to exact vocabulary: a word not seen during training contributes nothing to the prediction.

Fine-tuned BERT and its variants have dominated intent classification benchmarks since 2019 \citep{devlin2019bert}. The self-attention mechanism allows these models to capture contextual relationships between words, which helps with short, ambiguous utterances where the surrounding context determines intent. \citet{yamashita2023bert} showed that augmenting BERT with named entity recognition features improves performance on imbalanced call-center datasets.

More recently, decoder-based LLMs fine-tuned with parameter-efficient methods have been applied to classification tasks. \citet{sun2023text} showed that instruction-tuned LLMs can match or exceed BERT on standard classification benchmarks when fine-tuned with LoRA. However, their noise robustness has not been systematically compared to smaller models.

\subsection{ASR noise effects on NLU}

The impact of ASR errors on natural language understanding has been studied primarily in the context of encoder models. \citet{faruqui2021revisiting} showed that even moderate WER significantly reduces BERT's intent classification accuracy. \citet{shon2022slue} demonstrated that the gap between text-based and speech-based evaluation is domain-dependent and often larger than expected. Studies on the Audio-Snips dataset specifically quantified the BERT accuracy drop at multiple WER levels, providing the closest published baseline to our work.

A consistent finding across these studies is that WER is not a perfect predictor of downstream NLU degradation. The type of error matters: substitutions of content words (nouns, verbs) are more harmful than substitutions of function words \citep{faruqui2021revisiting}. Our noise simulation reflects this by weighting substitutions toward common English words rather than purely random vocabulary.

\subsection{The SLURP dataset}

\citet{bastianelli2020slurp} introduced SLURP as a large-scale spoken language understanding benchmark with 72,000 utterances across 18 domains and 66 intent classes. The dataset was collected in realistic home and office environments with variation in background noise, microphone quality, and speaker accent. The original paper provided baselines using earlier-generation ASR and NLU systems. No published work has applied QLoRA fine-tuned LLMs to this dataset, making it an ideal benchmark for our study. Our results constitute the first modern LLM classification baseline on SLURP.

\section{Methodology}

\subsection{Dataset}

We use the SLURP dataset accessed via HuggingFace (\texttt{qmeeus/slurp}), using the official train, development, and test splits: 50,342 training utterances, 8,649 validation utterances, and 13,006 test utterances across 66 intent classes. The dataset is moderately imbalanced: the largest class contains 3,247 training examples and the smallest contains fewer than 50. We address this during LLM training via capped minority class oversampling (see Section 3.3).

Table~\ref{tab:dataset} summarizes the dataset statistics relevant to our study.

\begin{table}[H]
\centering
\caption{SLURP dataset statistics used in this study.}
\label{tab:dataset}
\begin{tabular}{lc}
\toprule
\textbf{Property} & \textbf{Value} \\
\midrule
Total utterances & 72,277 \\
Training set & 50,342 \\
Validation set & 8,649 \\
Test set & 13,006 \\
Number of intent classes & 66 \\
Number of domains & 18 \\
Average utterance length & 6.2 words \\
Majority class (\% of test) & 4.97\% \\
Majority class baseline (Macro F1) & 0.0095 \\
\bottomrule
\end{tabular}
\end{table}

\subsection{ASR noise simulation}

We simulate ASR transcription errors at WER levels of 0, 10, 20, 30, 40, and 50\%. For each level, we generate a fixed noisy version of the test set using a deterministic random seed, so results are fully reproducible without running an actual ASR system.

The noise injection procedure follows the error type distribution reported in the ASR literature: 60\% substitutions, 25\% deletions, and 15\% insertions \citep{wang2003word}. This ratio reflects typical ASR behavior where substitutions are the most common error type.

Substitutions are performed using a hand-curated dictionary of 40 common English phonetic confusion pairs (for example, ``set'' substituted with ``get'', ``alarm'' with ``arm'', ``to'' with ``too''). For words not in the dictionary, character-level perturbations are applied: random adjacent character swaps, single character deletions, or single character substitutions with common English characters. Words shorter than four characters are replaced with random vocabulary items. Deletions remove randomly selected words from the utterance. Insertions add words sampled from the training vocabulary.

Table~\ref{tab:noise_example} in Appendix~\ref{app:noise} shows example noise injection output at each WER level. All noisy test sets are released as CSV files alongside this paper.

\subsection{Models}

\paragraph{TF-IDF + Linear SVM.}
We fit a TF-IDF vectorizer with 1-3 gram features, a maximum vocabulary of 50,000 terms, and sublinear TF scaling on the clean training set. A LinearSVC classifier with class-balanced weights is trained on the resulting feature matrix. This model represents the traditional machine learning approach that remains in widespread deployment.

\paragraph{TF-IDF + Logistic Regression.}
The same TF-IDF features with an L-BFGS logistic regression classifier ($C=5.0$, class-balanced weights, maximum 2,000 iterations). We include both TF-IDF variants because they represent the dominant traditional approach and show different noise sensitivity profiles.

\paragraph{BERT-base-uncased.}
We fine-tune BERT-base-uncased \citep{devlin2019bert} for 66-class sequence classification using the HuggingFace Trainer API. Training uses a learning rate of $3 \times 10^{-5}$, batch size 32, maximum sequence length 64 (appropriate for SLURP's short utterances), and up to 5 epochs with early stopping on validation Macro F1. BERT serves as the primary encoder baseline and represents the dominant NLP approach of the 2019-2022 period.

\paragraph{Qwen2.5-3B-Instruct (QLoRA).}
We fine-tune Qwen2.5-3B-Instruct \citep{qwen2024} using QLoRA \citep{dettmers2023qlora} with 4-bit NF4 quantization and bfloat16 compute. LoRA adapters are applied to all attention and feed-forward projection matrices ($r=32$, $\alpha=64$, dropout 0.1), yielding 59.9M trainable parameters out of 3.15B total (1.9\%).

We frame classification as a generative task. The system prompt lists all 66 intent class names and instructs the model to output the exact intent name as a single string. Training uses causal language modeling loss with all prompt tokens masked so that only the label name tokens contribute to the loss. Left-side truncation ensures the label token is never cut off during tokenization. The full class list is retained in the prompt at inference time to allow generalization to inputs the model has not seen during training.

Training hyperparameters: learning rate $5 \times 10^{-5}$ with cosine decay, effective batch size 16 (batch size 4, gradient accumulation 4), maximum sequence length 550, 3 epochs with early stopping patience 2. Minority class oversampling caps all classes at twice the median class count, yielding 80,621 training examples. All experiments were run on an NVIDIA RTX 4070 Super (12\,GB VRAM).

At inference, the model generates up to 12 tokens with greedy decoding. The generated text is matched to the nearest class name using exact match, then substring match.

\subsection{Evaluation}

We report Accuracy and Macro F1 as primary metrics. Macro F1 is the primary metric because it treats all 66 classes equally regardless of frequency, which is the appropriate measure when class imbalance is present and all intents are equally important in deployment. We also report Weighted F1 for completeness.

For the noise robustness analysis, we report the absolute Macro F1 drop from WER=0\% to WER=50\% as $\Delta$, and plot Macro F1 against WER level for each model.

All models are trained on clean text only. The evaluation protocol keeps all models on equal footing: each model sees clean training data and is evaluated on the same six noisy test sets.

\section{Results}

\subsection{Clean text performance}

Table~\ref{tab:clean_results} shows model performance on the clean (WER=0\%) test set.

\begin{table}[H]
\centering
\caption{Model performance on the clean SLURP test set (WER=0\%). Macro F1 is the primary metric.}
\label{tab:clean_results}
\begin{tabular}{lccc}
\toprule
\textbf{Model} & \textbf{Accuracy} & \textbf{Macro F1} & \textbf{Weighted F1} \\
\midrule
Majority class baseline & 0.0630 & 0.0095 & 0.0039 \\
TF-IDF + Linear SVM & 0.8260 & 0.7202 & 0.8149 \\
TF-IDF + Logistic Regression & 0.8305 & 0.7431 & 0.8221 \\
BERT-base (fine-tuned) & 0.8805 & 0.7896 & 0.8794 \\
Qwen2.5-3B QLoRA (ours) & \textbf{0.8843} & \textbf{0.7964} & \textbf{0.8832} \\
\bottomrule
\end{tabular}
\end{table}

The fine-tuned LLM achieves the highest Macro F1 of 0.7964, marginally ahead of BERT at 0.7896. The gap of 0.0068 is small but consistent across multiple runs. TF-IDF + LR trails at 0.7431 and TF-IDF + SVM at 0.7202, a gap of approximately 5-7 points behind the transformer models. This is consistent with the expectation that contextual representations provide some advantage on short, ambiguous utterances where surrounding context determines intent.

\subsection{Noise robustness}

Table~\ref{tab:noise_results} presents Macro F1 at all WER levels for each model. Figure~\ref{fig:main_plot} shows the degradation curves.

\begin{table}[H]
\centering
\caption{Macro F1 scores across WER levels on the SLURP test set (66 intent classes). $\Delta$ is the absolute drop from WER=0\% to WER=50\%.}
\label{tab:noise_results}
\begin{tabular}{lccccccc}
\toprule
\textbf{Model} & \textbf{0\%} & \textbf{10\%} & \textbf{20\%} & \textbf{30\%} & \textbf{40\%} & \textbf{50\%} & $\boldsymbol{\Delta}$ \\
\midrule
TF-IDF + SVM & 0.7202 & 0.6283 & 0.6323 & 0.6138 & 0.5848 & 0.5480 & 0.1722 \\
TF-IDF + LR  & 0.7431 & 0.6294 & 0.6345 & 0.6173 & 0.5977 & 0.5604 & 0.1827 \\
BERT-base    & 0.7896 & 0.6539 & 0.6422 & 0.6146 & 0.5745 & 0.5144 & \textbf{0.2752} \\
Qwen2.5-3B QLoRA & \textbf{0.7964} & \textbf{0.7079} & \textbf{0.7050} & \textbf{0.6737} & \textbf{0.6565} & \textbf{0.5888} & 0.2076 \\
\bottomrule
\end{tabular}
\end{table}

\begin{figure}[H]
\centering
\includegraphics[width=0.88\textwidth]{figures/paper1_fig1_macro_f1_vs_wer.png}
\caption{Macro F1 versus ASR word error rate for all four models on the SLURP test set. The LLM (Qwen2.5-3B QLoRA) maintains the highest Macro F1 at every noise level. BERT degrades most steeply despite comparable clean performance. TF-IDF models show a front-loaded degradation pattern.}
\label{fig:main_plot}
\end{figure}

Three findings are worth attention.

\textbf{Finding 1: The LLM maintains the highest Macro F1 at every noise level.} At WER=10\%, representing a high-quality commercial ASR system, the LLM leads BERT by 5.4 points (0.7079 vs 0.6539). This gap widens to 7.4 points at WER=50\% (0.5888 vs 0.5144). The advantage is consistent and grows with noise, suggesting the LLM has learned more noise-tolerant representations.

\textbf{Finding 2: BERT degrades most steeply despite competitive clean performance.} BERT's absolute drop of 28.2 points from clean to WER=50\% is the largest of all models. This is despite BERT having nearly the same clean Macro F1 as the LLM (0.7896 vs 0.7964). The contextual representations that give BERT its clean-text advantage appear to be more sensitive to token-level corruption than either bag-of-words features or the generative LLM's representations.

\textbf{Finding 3: TF-IDF models show an unexpected two-phase degradation.} Both TF-IDF variants drop sharply at WER=10\% (approximately 11-15 points) but then remain relatively stable from WER=10\% to WER=30\% before declining again. This plateau pattern differs qualitatively from the smooth monotonic decline of the transformer models. We hypothesize this occurs because TF-IDF depends on exact word matches, and once the high-frequency intent-discriminative keywords are corrupted at low WER, further substitutions affect lower-weight terms with diminishing impact on the classification score. The first hit is the most damaging.

Figure~\ref{fig:relative_plot} shows relative performance, normalizing each model to its WER=0\% baseline. This view highlights that BERT retains the smallest fraction of its clean performance at high noise levels, dropping to 65\% of baseline at WER=50\%, while the LLM retains 74\%.

\begin{figure}[H]
\centering
\includegraphics[width=0.88\textwidth]{figures/paper1_fig2_relative_degradation.png}
\caption{Relative Macro F1 as a percentage of each model's clean (WER=0\%) performance. BERT retains only 65\% of its baseline at WER=50\%, while the LLM retains 74\%.}
\label{fig:relative_plot}
\end{figure}

\subsection{Per-class analysis}

Looking at per-class results on the clean test set (see Appendix~\ref{app:per_class}), several patterns emerge across all models. High-frequency, semantically distinct classes such as \texttt{alarm\_remove}, \\ \texttt{transport\_taxi}, and \texttt{general\_greet} achieve near-perfect F1 across all models. The most difficult classes are those with semantic overlap or very few test examples: \texttt{general\_quirky} (LLM F1 = 0.66), \texttt{music\_settings} (0.51), and a handful of classes with fewer than 10 test examples that all models consistently misclassify.

The \texttt{general\_quirky} class is worth specific mention because it appears in the hardest-hit set for all models under noise. This class captures semantically ambiguous utterances that do not fit cleanly into any domain, and it is the class where LLM and BERT performance diverges most visibly at high WER. At WER=50\%, per-class analysis shows the LLM retains roughly 0.45 F1 on this class while BERT drops below 0.30. This suggests the LLM's broader pretraining allows it to handle unusual phrasing better when surface tokens are corrupted.

\section{Discussion}

\subsection{Practical implications}

Our results provide concrete, actionable guidance for practitioners building voice-based classification systems.

In low-noise environments (WER below 10\%), which represent well-controlled recording conditions or high-quality ASR, all four models achieve reasonable Macro F1 and the performance differences are modest. The choice between BERT and an LLM in this regime involves a trade-off between marginal accuracy gain and substantially higher inference cost and latency.

In high-noise environments (WER above 30\%), which represent telephone audio, accented speech, or challenging acoustic conditions, the performance gap becomes practically significant. At WER=40\%, the LLM leads BERT by 8.2 points (0.6565 vs 0.5745). At WER=50\%, the gap is 7.4 points. For a system processing thousands of calls per day, this difference translates directly into thousands of misclassifications. In this regime, the higher deployment cost of the LLM is more likely to be justified.

TF-IDF models occupy an interesting middle ground. Their absolute performance is lower than both transformer models across all noise levels, but their total degradation (17-18 points) is considerably less than BERT (28 points). A practitioner who cannot afford the compute cost of a transformer model may find that a well-tuned TF-IDF classifier provides acceptable and reasonably stable performance across noise conditions.

\subsection{Why LLMs are more robust}

We offer two hypotheses for the LLM's robustness advantage. First, scale: a 3B parameter model trained on orders of magnitude more text has encountered far more lexical variation, informal spelling, and noisy text during pretraining than BERT. This gives it inherently more robust representations for words and partial words that appear corrupted. Second, subword tokenization: while both BERT and Qwen use subword tokenization, Qwen's BPE vocabulary (around 152,000 tokens vs BERT's 30,000) is much larger, meaning it can represent corrupted words as meaningful subword units rather than falling back to an unknown token. A word corrupted from ``alarm'' to ``alrm'' is tokenized differently in Qwen, but the subword units still carry partial semantic information.

\subsection{Limitations}

We note three limitations. First, our noise simulation is synthetic. Real ASR errors have phonetic and language model biases our random perturbation procedure does not fully capture. Validating these findings against real ASR outputs at multiple WER levels is an important direction for future work. Second, we study a single dataset and a single LLM architecture and size. Generalizability to other domains, longer utterances, or other model families cannot be assumed. Third, the zero or near-zero F1 on several very small classes (fewer than 10 test examples) reflects a dataset limitation rather than a model limitation, and inflates the apparent variance in Macro F1 across noise levels for those classes.

\section{Conclusion}

We presented the first systematic comparison of TF-IDF, BERT, and a QLoRA fine-tuned LLM across controlled ASR noise levels on the SLURP spoken language understanding benchmark. Our main finding is that the fine-tuned LLM maintains higher Macro F1 at every noise level above zero, with a total degradation of 20.8 percentage points from clean to WER=50\%, compared to 28.2 points for BERT-base. This 7.4-point robustness advantage at WER=50\% is practically significant for real-world voice applications. We also identified a previously unreported two-phase degradation pattern in TF-IDF classifiers, and we established the first published LLM classification baselines on the SLURP dataset.

Our findings suggest that for deployments where ASR quality cannot be guaranteed, fine-tuned LLMs represent a more robust choice than BERT despite higher inference cost. Future work should validate these findings with real ASR outputs and extend the comparison to larger LLMs and additional intent classification benchmarks.

\bibliographystyle{plainnat}
\bibliography{references}

\appendix

\section{Noise injection examples}
\label{app:noise}

\begin{table}[H]
\centering
\caption{Example noise injection output at each WER level for a single test utterance.}
\label{tab:noise_example}
\begin{tabular}{cp{10cm}}
\toprule
\textbf{WER} & \textbf{Utterance} \\
\midrule
0\%  & set an alarm for seven in the morning \\
10\% & set league alarm for seven in the morning \\
20\% & set league alarm for alan seven in the morning \\
30\% & set league alarm for alan seven in the morning \\
40\% & get following for seven in the morning \\
50\% & get convenient following for seven in the morning \\
\bottomrule
\end{tabular}
\end{table}

\section{Per-class classification report}
\label{app:per_class}

The full per-class precision, recall, and F1 scores for the LLM at WER=0\% are provided in the supplementary file \texttt{llm\_report\_wer0.txt}, available in the code repository. Selected highlights for the hardest and easiest classes are shown below.

\textit{Easiest classes (LLM F1 at WER=0\%):} \texttt{alarm\_remove} (1.000), \texttt{transport\_taxi} (1.000), \texttt{general\_greet} (1.000), \texttt{alarm\_query} (0.980), \texttt{play\_game} (0.992).

\textit{Hardest classes (LLM F1 at WER=0\%):} \texttt{general\_quirky} (0.664), \texttt{music\_settings} (0.512), \texttt{iot\_hue\_lighton} (0.182), \texttt{joke} (0.417), \texttt{podcasts} (0.200).

Zero-F1 classes (\texttt{factoid}, \texttt{music}, \texttt{query}, \texttt{quirky}, \texttt{remove}) each have fewer than 25 test examples with no training signal sufficient for any model to learn them. This is a known dataset limitation in SLURP and does not reflect model failure.

\section{Implementation details}
\label{app:impl}

All experiments were run on Windows 11 with an NVIDIA RTX 4070 Super (12\,GB VRAM). Software: Python 3.11, PyTorch 2.11.0, HuggingFace Transformers 5.4.0, PEFT, bitsandbytes. BERT training took approximately 35 minutes. LLM training took approximately 9 hours. Noise evaluation took approximately 2 hours for TF-IDF models and BERT, and approximately 2 hours per WER level for the LLM.

All code, model adapter weights, and noisy test CSV files are available at: \\\texttt{https://github.com/M-O-J-U/SLURP}.

\end{document}
